% GENERATED FILE. Edit canonical lessons, then run npm run generate:books.
% canonical-source-hash: fnv1a64:9bf0040a99874fc7
% canonical-lessons: IT-C456-parrucchiere, IT-C456-barbiere, IT-C456-estetista, IT-C456-ottico, IT-C456-smacchiare

\chapter{Capelli e occhiali}
\label{ch:it-capelli-e-occhiali}
\hlchaptermodality{456}

\begin{chapteropening}
\textbf{By the end of this chapter:} \emph{I can use personal services: the hairdresser, the barber, the beautician, the optician and getting a stain out.}

Complete the last lesson of chapter 456: I can use personal services: the hairdresser, the barber, the beautician, the optician and getting a stain out.
\end{chapteropening}

\section[il parrucchiere]{il parrucchiere — the hairdresser}

\label{lesson:IT-C456-parrucchiere}



\begin{quote}
Before the new one: say the Italian for to shorten, then the Italian for the locksmith.
\end{quote}

\subsection*{You'll want to know: il parrucchiere}
\textbf{il parrucchiere} — \textquotedblleft{}the hairdresser\textquotedblright{}.

\begin{quote}
Ho appuntamento dal parrucchiere alle tre.
\end{quote}

\begin{quote}
\emph{I have an appointment at the hairdresser's at three.}
\end{quote}

\emph{La parrucchiera}; \emph{andare dal parrucchiere}.

\begin{cousinweb}
\emph{parrucchiere} is built on \emph{parrucca} 'wig': first the wig-maker. The origin of \emph{parrucca} is not settled.
\end{cousinweb}

\begin{grammarlens}[title={What you've built}]
One more word for telling, explaining and describing at B1.
\end{grammarlens}

\subsection*{Your turn}
\begin{itemize}
\raggedright
  \item \emph{Say it:} \emph{il parrucchiere}
  \item \emph{Say it:} \emph{il parrucchiere}, once more
  \item \emph{Say it:} a sentence of your own with \emph{il parrucchiere}
  \item \emph{Recall:} say the Italian for to shorten, then the Italian for the locksmith, then say \emph{il parrucchiere} again
\end{itemize}

\begin{tcolorbox}[breakable,colback=teal!4,colframe=teal!35!black,title={Before you move on}]
What does il parrucchiere mean? (\textquotedblleft{}The hairdresser\textquotedblright{}.) Say it once more.
\end{tcolorbox}
\section[il barbiere]{il barbiere — the barber}

\label{lesson:IT-C456-barbiere}



\begin{quote}
Before the new one: say the Italian for the locksmith, then the Italian for the hairdresser.
\end{quote}

\subsection*{You'll want to know: il barbiere}
\textbf{il barbiere} — \textquotedblleft{}the barber\textquotedblright{}.

\begin{quote}
Mio nonno va dal barbiere ogni sabato.
\end{quote}

\begin{quote}
\emph{My grandfather goes to the barber every Saturday.}
\end{quote}

\emph{Farsi la barba} — to shave.

\begin{cousinweb}
\emph{barbiere} is built on \emph{barba} 'beard' (EtymDB), which is inherited from Latin \emph{barba}.
\end{cousinweb}

\begin{grammarlens}[title={What you've built}]
One more word for telling, explaining and describing at B1.
\end{grammarlens}

\subsection*{Your turn}
\begin{itemize}
\raggedright
  \item \emph{Say it:} \emph{il barbiere}
  \item \emph{Say it:} \emph{il barbiere}, once more
  \item \emph{Say it:} a sentence of your own with \emph{il barbiere}
  \item \emph{Recall:} say the Italian for the locksmith, then the Italian for the hairdresser, then say \emph{il barbiere} again
\end{itemize}

\begin{tcolorbox}[breakable,colback=teal!4,colframe=teal!35!black,title={Before you move on}]
What does il barbiere mean? (\textquotedblleft{}The barber\textquotedblright{}.) Say it once more.
\end{tcolorbox}
\section[l'estetista]{l'estetista — the beautician}

\label{lesson:IT-C456-estetista}



\begin{quote}
Before the new one: say the Italian for the hairdresser, then the Italian for the barber.
\end{quote}

\subsection*{You'll want to know: l'estetista}
\textbf{l'estetista} — \textquotedblleft{}the beautician\textquotedblright{}.

\begin{quote}
L'estetista lavora vicino al parrucchiere.
\end{quote}

\begin{quote}
\emph{The beautician works near the hairdresser's.}
\end{quote}

One form for a man or a woman.

\begin{cousinweb}
\emph{estetista} is built on \emph{estetica} 'beauty care', a learned borrowing from Greek \emph{aisthētikós} 'of the senses'.
\end{cousinweb}

\begin{grammarlens}[title={What you've built}]
One more word for telling, explaining and describing at B1.
\end{grammarlens}

\subsection*{Your turn}
\begin{itemize}
\raggedright
  \item \emph{Say it:} \emph{l'estetista}
  \item \emph{Say it:} \emph{l'estetista}, once more
  \item \emph{Say it:} a sentence of your own with \emph{l'estetista}
  \item \emph{Recall:} say the Italian for the hairdresser, then the Italian for the barber, then say \emph{l'estetista} again
\end{itemize}

\begin{tcolorbox}[breakable,colback=teal!4,colframe=teal!35!black,title={Before you move on}]
What does l'estetista mean? (\textquotedblleft{}The beautician\textquotedblright{}.) Say it once more.
\end{tcolorbox}
\section[l'ottico]{l'ottico — the optician}

\label{lesson:IT-C456-ottico}



\begin{quote}
Before the new one: say the Italian for the barber, then the Italian for the beautician.
\end{quote}

\subsection*{You'll want to know: l'ottico}
\textbf{l'ottico} — \textquotedblleft{}the optician\textquotedblright{}.

\begin{quote}
Dall'ottico ho scelto occhiali nuovi.
\end{quote}

\begin{quote}
\emph{At the optician's I chose new glasses.}
\end{quote}

\emph{Dall'ottico} — at the optician's.

\begin{cousinweb}
\emph{ottico} is a learned borrowing from Greek \emph{optikós} 'of sight', through Medieval Latin \emph{opticus}.
\end{cousinweb}

\begin{grammarlens}[title={What you've built}]
One more word for telling, explaining and describing at B1.
\end{grammarlens}

\subsection*{Your turn}
\begin{itemize}
\raggedright
  \item \emph{Say it:} \emph{l'ottico}
  \item \emph{Say it:} \emph{l'ottico}, once more
  \item \emph{Say it:} a sentence of your own with \emph{l'ottico}
  \item \emph{Recall:} say the Italian for the barber, then the Italian for the beautician, then say \emph{l'ottico} again
\end{itemize}

\begin{tcolorbox}[breakable,colback=teal!4,colframe=teal!35!black,title={Before you move on}]
What does l'ottico mean? (\textquotedblleft{}The optician\textquotedblright{}.) Say it once more.
\end{tcolorbox}
\section[smacchiare]{smacchiare — to remove stains from}

\label{lesson:IT-C456-smacchiare}



\begin{quote}
Before the new one: say the Italian for the beautician, then the Italian for the optician.
\end{quote}

\subsection*{You'll want to know: smacchiare}
\textbf{smacchiare} — \textquotedblleft{}to remove stains from\textquotedblright{}.

\begin{quote}
In lavanderia possono smacchiare la giacca.
\end{quote}

\begin{quote}
\emph{At the dry cleaner's they can get the stains out of the jacket.}
\end{quote}

\emph{La macchia} — a stain.

\begin{cousinweb}
\emph{smacchiare} is built from \emph{s-} and \emph{macchia} 'stain', which is inherited from Latin \emph{macula} 'spot' (Lewis \& Short).
\end{cousinweb}

\begin{grammarlens}[title={What you've built}]
One more word for telling, explaining and describing at B1.
\end{grammarlens}

\subsection*{Your turn}
\begin{itemize}
\raggedright
  \item \emph{Say it:} \emph{smacchiare}
  \item \emph{Say it:} \emph{smacchiare}, once more
  \item \emph{Say it:} a sentence of your own with \emph{smacchiare}
  \item \emph{Recall:} say the Italian for the beautician, then the Italian for the optician, then say \emph{smacchiare} again
\end{itemize}

\begin{tcolorbox}[breakable,colback=teal!4,colframe=teal!35!black,title={Before you move on}]
What does smacchiare mean? (\textquotedblleft{}To remove stains from\textquotedblright{}.) Say it once more.
\end{tcolorbox}
